En el colegio, los estudiantes de grado noveno participan en un proyecto de servicio en el que deben pintar las columnas del patio principal. Una de estas columnas tiene forma de prisma hexagonal regular.

Cada lado de la base mide 40 cm y la altura de la columna es de 3 m. Los estudiantes saben que no es necesario pintar la parte superior ni inferior de la columna, ya que estas no están expuestas.

Además, el proveedor indica que se requieren aproximadamente 0,25 litros de pintura por cada metro cuadrado de superficie.

Con base en esta información, determine cuánta pintura se debe comprar para pintar completamente la columna.

Explique su procedimiento, las decisiones tomadas y los conceptos geométricos utilizados.